% ============================================================
% Chapter 5 — Results (target 14–16 pp)
% Written LAST among the big chapters, once the 2x2 campaign (n=1 seed,
% decided 2026-07-14: S2/S3 dropped for compute budget) is consolidated.
% Every number recomputed from episodes.jsonl on disk.
% Data status (2026-07-14 evening): MLP-C 20260622_171626 train+eval OK;
% MLP-B 20260623_171855 train+eval OK; GNN-C 20260630_181143 train only;
% GNN-B 20260714_155209 training in flight; final evals remain for the
% two GNN runs only.
% ============================================================
\chapter{Results}
\label{ch:results}

\todo{Awaiting campaign consolidation. Section skeleton below.}

% Sources: recomputed from disk 2026-07-14 (9 runs × episodes.jsonl,
% dedup + route_complete only; integrity 6/6, 0 bad lines on all).
% Figure: docs/thesis/scripts/make_pilot_collapse_fig.py.
\section{Preliminary Architectural Pilot (April 2026)}
\label{sec:pilot}

Before the correctness fixes of \cref{sec:bugs}, nine runs (April--early
May 2026) crossed critic encoders and training regimes in the
first-generation environment: seed 42, budgets of $1$M or $3$M steps,
44D-era vehicle observations, and the pre-fix route planner. They are
reported here as a preliminary study --- not because their comparisons
hold (they do not), but because they motivated the methodology of
\cref{ch:methodology} and provide this thesis's clearest cautionary
example about aggregate metrics. All nine logs pass the integrity checks
of \cref{sec:metrics}; \cref{tab:pilot} reports the recomputed
cumulative metrics.

\begin{table}[t]
  \centering\small
  \caption{April 2026 pilot runs, recomputed from disk (cumulative,
  agent-level; \%). All runs: seed 42, Town03, pre-route-length-fix
  environment. Per-record speed was not yet logged in this era.
  $^{a}$~architecture untagged in run metadata (old configuration
  schema); confirmed MLP by the author. $^{b}$~variant with attention critic encoder
  and PopArt enabled. $^{c}$~stopped at $\approx$2.7M of 3M.}
  \label{tab:pilot}
  \begin{tabular}{@{}llllrrrrrr@{}}
    \toprule
    Run & Regime & Critic & Steps & Ep. & \multicolumn{4}{c}{Vehicles} & Ped. \\
    \cmidrule(lr){6-9}\cmidrule(l){10-10}
    & & & & & SR & s$+$t & coll. & off. & SR \\
    \midrule
    0406 & curr.\ & MLP & 1M & 1062 & 1.2 & 46.4 & 31.8 & 20.7 & 82.4 \\
    0411 & curr.\ & MLP & 3M & 3047 & 2.0 & 53.2 & 30.6 & 14.2 & 67.4 \\
    0425 & curr.\ & GNN & 3M & 3015 & 1.2 & 70.8 & 19.5 & 8.6 & 70.6 \\
    0430 & curr.\ & GNN$^{b}$ & 3M & 2188 & 1.8 & 50.9 & 28.9 & 18.5 & 66.2 \\
    0504 & curr.\ & GNN & 3M$^{c}$ & 2730 & 3.9 & 64.3 & 21.1 & 10.8 & 75.4 \\
    0405 & batch & MLP & 1M & 1066 & 1.9 & 42.3 & 38.3 & 17.6 & 77.0 \\
    0408 & batch & MLP$^{a}$ & 3M & 3038 & 0.2 & 57.8 & 28.2 & 13.8 & 72.5 \\
    0421 & batch & GNN & 3M & 3078 & 2.7 & 49.5 & 30.5 & 17.4 & 74.4 \\
    0502 & batch & GNN & 3M & 3054 & 3.0 & 58.2 & 23.6 & 15.3 & 74.4 \\
    \bottomrule
  \end{tabular}
\end{table}

Four observations, in decreasing order of certainty:

\begin{enumerate}
  \item \textbf{Vehicles sat at floor in every cell} (SR 0.2--3.9\%).
        At floor, the success metric has no room to express differences:
        whatever effect architecture or regime may have had is compressed
        below the noise. The pilot's architectural comparison on vehicles
        is therefore \emph{uninformative} --- the counterexample that
        motivates running the main campaign in the mid-range regime. The
        floor itself was later traced to the route-length defect of
        \cref{sec:bugs}.
  \item \textbf{Pedestrians learned} (SR 66--82\%): the platform was
        trainable for the simpler class even in this era, localizing the
        failure to the vehicle task rather than to the stack --- and
        illustrating why all reporting is per class
        (\cref{sec:metrics}).
  \item \textbf{A qualitative architectural signature.} Among the
        curriculum runs, the plain GNN-critic run (0425) shows the most
        defensive failure profile of the nine: the highest
        stuck$+$timeout (70.8\%) with the lowest collision (19.5\%) and
        offroad (8.6\%), versus 53.2/30.6/14.2\% for its MLP counterpart
        (0411). At floor SR this is hypothesis-generating at most;
        whether a regime-dependent architectural signature reappears on
        the healthy trunk is exactly question H3, taken up in
        \cref{sec:arch-axis}.
  \item \textbf{Cumulative aggregates can hide a collapse.} Run 0504 was
        flagged during training as having collapsed mid-run --- yet its
        \emph{cumulative} vehicle SR (3.9\%) is the highest of all nine.
        \Cref{fig:pilot-collapse} resolves the paradox: the run sustained
        a windowed SR of 10--14\% for roughly 500 episodes, then dropped
        abruptly to near zero around episode 1{,}100 and never recovered;
        the cumulative figure is entirely inherited from the
        pre-collapse phase. This lesson --- current competence must be
        measured on a window, not since birth --- directly shaped the
        windowed unlock metric of \cref{sec:regimes} and the quartile
        reporting of \cref{sec:metrics}.
\end{enumerate}

\begin{figure}[t]
  \centering
  \includegraphics[width=0.85\textwidth]{pilot_collapse_timeseries}
  \caption{Windowed vehicle SR (100-episode window) for the two plain
  GNN-curriculum pilot runs. Run 0504 (dashed) sustains 10--14\% for
  $\approx$500 episodes, collapses near episode 1{,}100, and never
  recovers --- while ending with the \emph{highest} cumulative SR of the
  pilot (3.9\%, \cref{tab:pilot}). Cumulative aggregates inherit
  history; collapse detection requires windowed views.}
  \label{fig:pilot-collapse}
\end{figure}

% Sources for §5.2–5.4, §5.6: docs/thesis/data/campaign_consolidation.json,
% recomputed from disk 2026-07-15 by docs/thesis/scripts/
% consolidate_campaign.py (dedup, route_complete only, per class).
% GNN-B cells marked \todo until its training + the two GNN evals land.
\section{Campaign Integrity and Setup Recap}
\label{sec:integrity}

\Cref{tab:runmatrix} lists the campaign runs. All were launched on the
same trunk (commit \texttt{7defb03}) with seed 999; configuration diffs
confirm the cells differ only in the regime and critic-encoder flags
(\cref{sec:campaign}). The three completed training logs pass every
integrity condition of \cref{sec:metrics}: exactly six records per
episode throughout, zero malformed lines, zero non-finite values.

\begin{table}[t]
  \centering\small
  \caption{Campaign run matrix (all: 3M steps, Town03, seed 999, trunk
  \texttt{7defb03}). Realized level exposure is reported as
  easy/medium/hard shares of episodes.}
  \label{tab:runmatrix}
  \begin{tabular}{@{}llrrlll@{}}
    \toprule
    Cell & Run ID & Episodes & Records & Level mix (\%) & Integrity & Final eval \\
    \midrule
    MLP-curr.\ & 20260622\_171626 & 3220 & 19{,}320 & 32.5/45.8/21.6 & OK & 400 ep, done \\
    MLP-batch  & 20260623\_171855 & 3180 & 19{,}080 & 35.3/32.8/31.9 & OK & 400 ep, done \\
    GNN-curr.\ & 20260630\_181143 & 3293 & 19{,}758 & 32.4/37.4/30.1 & OK & \todo{pending} \\
    GNN-batch  & 20260714\_155209 & \multicolumn{4}{l}{\todo{training in progress}} & \todo{pending} \\
    \bottomrule
  \end{tabular}
\end{table}

The realized exposure matches the two designs: the batch run stays within
$\approx 2$ percentage points of the uniform $1/3$ per level, with all
three levels present within its first twenty episodes; the curriculum
runs concentrate easy/medium exposure early and reach hard only after
their unlocks (\cref{sec:sample-efficiency}), ending with hard shares of
21.6\% (MLP) and 30.1\% (GNN).

% ------------------------------------------------------------
\section{Sample Efficiency}
\label{sec:sample-efficiency}

\begin{figure}[t]
  \centering
  \includegraphics[width=\textwidth]{learning_curves}
  \caption{Windowed success rate (100-episode window) versus environment
  steps for the completed cells. \todo{Re-run
  \texttt{make\_fig5\_learning\_curves.py} to add GNN-batch when its
  training completes.} Important caveat: each curve is measured on that
  run's \emph{own concurrent level mix} --- gated for the curriculum
  cells, uniform for batch --- so vertical comparisons between arms are
  not like-for-like until the curriculum has unlocked all levels; the
  distribution-matched comparison is the final evaluation
  (\cref{sec:final-performance}).}
  \label{fig:learning-curves}
\end{figure}

\paragraph{Early phase: indistinguishable take-off.}
All three completed cells climb at the same speed: the 100-episode
windowed vehicle SR first crosses 30\% after 137--140K steps and 50\%
after 162--167K steps in every cell --- a spread of a few percent,
far inside noise. The notable part is \emph{what} each policy crossed it
on: the curriculum cells on easy-only exposure, the batch cell on the
full uniform mix including \SI{100}{m} routes. At equal budget, early
success-rate take-off was not the discriminating factor between regimes
in this campaign.

\paragraph{Unlock timing (curriculum cells).}
Medium unlocked at practically the same point in both curriculum runs
(episode 842 of 3220, 26.1\%, MLP; episode 860 of 3293, 26.1\%, GNN).
Hard, however, unlocked much earlier under the GNN critic: episode
1{,}910 (58.0\% of training) versus 2{,}286 (71.0\%) for MLP --- the
latter consistent with the 70\%-budget pressure cap firing rather than
the competence gate (\cref{sec:regimes}), while the GNN cell reached the
hard gate on competence. The result is visible in the exposure shares of
\cref{tab:runmatrix}: the GNN-curriculum policy saw $\approx 40$\% more
hard episodes than its MLP counterpart.

\paragraph{After the take-off.}
From $\approx 0.5$M steps the MLP-batch curve settles visibly below the
two curriculum curves (\cref{fig:learning-curves}), which track each
other closely; all curves drift downward late in training as the
curriculum mixes harden (curriculum cells) and as pedestrian windowed
success declines --- the within-run dynamics are analysed in
\cref{sec:behavior}.

% ------------------------------------------------------------
\section{Final Cumulative Performance}
\label{sec:final-performance}

\begin{table}[t]
  \centering\small
  \caption{Cumulative training metrics per cell (agent-level, \%;
  speed in km/h for vehicles, m/s for pedestrians).
  \todo{Add the GNN-batch row when its training completes.}}
  \label{tab:cumulative}
  \begin{tabular}{@{}lrrrrrrrr@{}}
    \toprule
    & \multicolumn{5}{c}{Vehicles} & \multicolumn{3}{c}{Pedestrians} \\
    \cmidrule(lr){2-6}\cmidrule(l){7-9}
    Cell & SR & s$+$t & coll. & off. & speed & SR & s$+$t & speed \\
    \midrule
    MLP-curr.\ & 57.28 & 30.48 & 8.84 & 3.41 & 10.2 & 56.59 & 42.56 & 1.96 \\
    MLP-batch  & 46.31 & 23.61 & 17.33 & 12.76 & 18.7 & 57.06 & 42.03 & 1.61 \\
    GNN-curr.\ & 56.55 & 20.84 & 13.37 & 9.23 & 18.0 & 60.50 & 38.58 & 1.91 \\
    GNN-batch  & \multicolumn{8}{l}{\todo{training in progress}} \\
    \bottomrule
  \end{tabular}
\end{table}

\begin{table}[t]
  \centering\small
  \caption{Final 400-episode evaluation, MLP pair (agent-level, \%).
  Scenarios: easy/medium/hard on Town03, test on Town05.
  \todo{Twin table for the GNN pair once its evaluations run.}}
  \label{tab:eval-mlp}
  \begin{tabular}{@{}llrrrrr@{}}
    \toprule
    Scenario & Cell & veh SR & veh s$+$t & veh coll. & veh off. & ped SR \\
    \midrule
    easy   & MLP-curr.\ & 60.33 & 28.00 & 7.33 & 4.33 & 58.00 \\
           & MLP-batch  & 41.33 & 28.33 & 18.00 & 12.33 & 58.67 \\
    medium & MLP-curr.\ & 51.67 & 36.00 & 9.00 & 3.33 & 40.67 \\
           & MLP-batch  & 31.00 & 27.67 & 25.33 & 16.00 & 41.67 \\
    hard   & MLP-curr.\ & 44.33 & 38.33 & 14.33 & 3.00 & 26.67 \\
           & MLP-batch  & 29.00 & 29.67 & 23.33 & 18.00 & 32.33 \\
    test   & MLP-curr.\ & 37.00 & 50.67 & 11.00 & 1.33 & 62.67 \\
           & MLP-batch  & 18.00 & 39.00 & 37.33 & 5.67 & 63.00 \\
    \midrule
    all    & MLP-curr.\ & 48.33 & 38.25 & 10.42 & 3.00 & 47.00 \\
           & MLP-batch  & 29.83 & 31.17 & 26.00 & 13.00 & 48.92 \\
    \bottomrule
  \end{tabular}
\end{table}

\paragraph{The primary contrast, MLP pair.}
On the distribution-matched final evaluation
(\cref{tab:eval-mlp}), the curriculum-trained vehicle policy completes
48.33\% of routes against 29.83\% for batch: $+18.5$ percentage points,
with the same sign in \emph{all four} scenarios ($+19.0$, $+20.7$,
$+15.3$, $+19.0$\,pp) and in cumulative training ($+11.0$\,pp,
\cref{tab:cumulative}). The margin is roughly two to four times the
same-seed replicate band of \cref{sec:iterations}, and its consistency
across scenarios, difficulty levels and measurement layers is what the
paired design can offer in place of seed replication
(\cref{sec:campaign}).

\paragraph{Two different ways to fail.}
The regimes produce recognizably different driving. The batch policy is
nearly twice as fast (18.7 vs.\ 10.2 km/h in training) but converts that
mobility into contact: collision 17.3\% vs.\ 8.8\% and offroad 12.8\%
vs.\ 3.4\% in training, widening to collision 26.0\% vs.\ 10.4\% at
evaluation. The curriculum policy fails mostly by stalling
(stuck$+$timeout 38.3\% vs.\ 31.2\% at evaluation). Neither profile
dominates the other on every axis --- the curriculum trades speed for
safety --- but on the strict success criterion the trade resolves
clearly in its favour.

\paragraph{Pedestrians: no regime effect.}
Pedestrian success is flat between arms everywhere: 56.6 vs.\ 57.1\% in
training, 47.0 vs.\ 48.9\% at evaluation, and within $\pm 6$\,pp in
every scenario. The training regime, in this platform, is a
vehicle-policy variable; the pedestrian task appears insensitive to task
ordering. Pedestrian failure structure is taken up in
\cref{sec:ped-limits}.

\paragraph{GNN-curriculum (training-only preview).}
Pending its evaluation, the GNN-curriculum cell lands within noise of
MLP-curriculum on headline SR (56.55 vs.\ 57.28\%) while redistributing
the failure mass: less stalling ($-9.6$\,pp stuck$+$timeout), more
contact ($+4.5$\,pp collision, $+5.8$\,pp offroad), and much higher
speed (18.0 vs.\ 10.2 km/h) --- closer to the batch cell's mobility with
curriculum-level success. If confirmed at evaluation, the critic encoder
would be shaping the \emph{risk profile} more than the success rate;
this thread is picked up in \cref{sec:arch-axis}.

\section{Behavioral Analysis}\label{sec:behavior}
% 2–2.5 pp: Q1->Q4 trajectories (data ready in campaign_consolidation.json:
% all cells decline Q1->Q4; MLP-B declines on a CONSTANT mix — real late
% degradation vs mix-hardening for C cells — disentangle), speed profiles,
% no_wp_steps, stuck causes; defensive-equilibrium case study (hazard gate)
% + P1xV1 coupling; ped windowed decline (fig) mix-confounded for C.

\section{Generalization: Town03 to Town05}
\label{sec:generalization}

The test scenario places the same policies on Town05 --- a squared-grid,
multi-lane layout structurally unlike Town03 (\cref{sec:carla-bg}) ---
for 100 episodes at \SI{80}{m} targets, with traffic as in training.
Three observations from the MLP pair (\cref{tab:eval-mlp}); the GNN pair
is \todo{pending evaluation}.

\paragraph{The ordering survives the map change.}
Curriculum completes 37.0\% of Town05 routes against batch's 18.0\%
($+19.0$\,pp) --- essentially the same margin as in-domain. The
absolute drop from the Town03 scenario average is nearly identical for
the two arms ($-15.1$\,pp from 52.1\% for curriculum; $-15.8$\,pp from
33.8\% for batch), so the curriculum advantage is not a Town03
idiosyncrasy: it transfers whole to roads never seen in training.
Read against hypothesis H2, the broader early coverage of the batch
regime bought it \emph{no} generalization edge --- if anything the
relative drop is worse from its lower base ($-47$\% vs.\ $-29$\%
relative).

\paragraph{Failure signatures amplify off-distribution.}
Each arm fails on Town05 the way it fails at home, only more so: the
curriculum policy stalls (stuck$+$timeout 50.7\%, its highest anywhere)
while keeping contact low (collision 11.0\%, offroad 1.3\%); the batch
policy crashes (collision 37.3\%, its highest anywhere). Novel roads
stress exactly the axis each policy already leaned on.

\paragraph{Pedestrians transfer easily.}
Pedestrian success on Town05 (62.7 vs.\ 63.0\%) exceeds both arms'
Town03 hard-level results by $\approx 30$\,pp --- the \SI{80}{m} test
routes are more feasible than Town03's \SI{100}{m} hard routes, pointing
again at route feasibility rather than policy generalization as the
pedestrian bottleneck (\cref{sec:ped-limits}).

\section{Architecture Axis: MLP vs GNN Critic}\label{sec:arch-axis}
% 1.5–2 pp (H3): does the C-vs-B contrast replicate on both architectures?
% arch x regime interaction; link to April pilot signature.

\section{Pedestrian Structural Limits}\label{sec:ped-limits}
% 1–1.5 pp: route feasibility (under_target/too_short/sidewalk_fallback),
% medium/hard collapse: environment vs policy limits.

\section{Findings versus Hypotheses}\label{sec:findings}
% 0.5–1 pp: H1/H2/H3 -> supported/rejected/partial per axis.
